\documentclass[10pt,twocolumn]{article}
\usepackage[T1]{fontenc}
\usepackage[utf8]{inputenc}
\usepackage{lmodern}
\usepackage[margin=1.8cm,top=2.4cm,bottom=2cm,columnsep=0.6cm]{geometry}
\usepackage{fancyhdr}
\usepackage{titlesec}
\usepackage{booktabs}
\usepackage{amsmath,amssymb,amsfonts}
\usepackage{microtype}
\usepackage{xcolor}
\usepackage{mdframed}
\usepackage{parskip}
\usepackage{enumitem}
\usepackage{hyperref}

\definecolor{rulered}{RGB}{180,30,80}
\definecolor{boxgray}{RGB}{245,245,245}
\definecolor{keywordgray}{RGB}{100,100,100}

\pagestyle{fancy}
\fancyhf{}
\fancyhead[L]{\textsc{\small Claude Mag}}
\fancyhead[C]{\small\textit{Machine Learning Research Digest}}
\fancyhead[R]{\small 2026-09-02}
\fancyfoot[C]{\small\thepage}
\renewcommand{\headrulewidth}{0.8pt}

\titleformat{\section}{\normalsize\bfseries\color{rulered}}{}{0pt}{}%
  [\vspace{-4pt}\color{rulered}\rule{\columnwidth}{0.4pt}]
\titlespacing{\section}{0pt}{8pt}{4pt}

\hypersetup{colorlinks=true,urlcolor=rulered,linkcolor=black}

\begin{document}
\setlength{\parindent}{0pt}
\setlength{\parskip}{4pt}

{\small\color{keywordgray}\textbf{Keywords:}
  vision-language models \textbullet{}
  assistive technology \textbullet{}
  salience \textbullet{}
  low vision}\par\vspace{4pt}

{\LARGE\bfseries\raggedright
  Focus Where It Counts: A Salience-Driven
  Vision-Language Model for Low Vision Assistance}\par\vspace{4pt}

{\large\itshape\raggedright
  Captioning That Prioritizes Scene Elements by Human-Centered
  Importance, Not Photographic Coverage}\par\vspace{6pt}

{\small\textbf{By} Jiazhao Liang, Hao Huang, Shuaihang Yuan, et al.}\par\vspace{2pt}

{\small\color{keywordgray}arXiv:2608.28218 \textbullet{} ECCV 2026}
\par\vspace{2pt}\noindent{\color{rulered}\rule{\columnwidth}{0.6pt}}\vspace{6pt}

Vision-language models offer powerful scene captioning but are designed
for comprehensive coverage, not for the ordered priorities of users with
blindness or low vision. This paper introduces Salience-LLaVA, a
salience-driven VLM trained to mention scene elements in order of their
functional importance for human-centered navigation and interaction.
Supported by three newly curated salience-annotated datasets, the model
places the highest-priority element first in 84.3\% of test images
and is rated 4.2/5 by blind and low-vision users—versus 3.1/5 for
standard LLaVA.

\begin{mdframed}[backgroundcolor=boxgray,linecolor=rulered,linewidth=1pt,
    innerleftmargin=6pt,innerrightmargin=6pt,
    innertopmargin=6pt,innerbottommargin=6pt]
\textbf{\small AT A GLANCE}\par\vspace{4pt}
\begin{description}[leftmargin=0pt,labelsep=4pt,itemsep=2pt,parsep=0pt,
    font=\small\bfseries]
  \item[Problem:] General-purpose VLMs describe images comprehensively
    but do not order information by functional relevance for users with
    blindness or low vision.
  \item[Why it matters:] BLV users often hear only the first few seconds
    of a caption; placing important information last renders AI assistance
    unreliable for navigation and safety.
  \item[Method:] Curate salience-annotated datasets; fine-tune LLaVA
    with a combined captioning and ranking loss to produce salienceordered
    descriptions; prepend salience scores as soft prompts.
  \item[Result:] 84.3\% top-1 salience ordering accuracy (vs.\ 51.2\%
    for LLaVA); 4.2/5 user helpfulness rating vs.\ 3.1/5.
  \item[Evidence:] Three datasets, user study with BLV participants,
    transfer experiments on VizWiz.
\end{description}
\end{mdframed}

\section{The Big Picture}

Assistive technology for persons with blindness or low vision (BLV) is
undergoing a rapid transformation as vision-language models replace
earlier rule-based systems. VLMs can describe complex scenes in natural
language, answer visual questions, and identify objects---capabilities
that are transformative for BLV users navigating the world.

But VLMs are trained on captioning benchmarks that reward comprehensive
coverage: a high CIDEr score requires mentioning many objects accurately.
For a BLV user in a grocery store, a caption that begins with the color
and pattern of the floor tiles before mentioning that the aisle is blocked
by a forklift is not just unhelpful---it is dangerous.

This paper identifies the ordering problem as a first-class challenge in
assistive VLM design and proposes a framework that addresses it through
salience-aware training.

\section{Why Existing Approaches Fall Short}

\textbf{General-purpose captioning:} Models trained on MSCOCO and
similar benchmarks describe images from top-left to bottom-right, or
by object size/saliency in the photographic sense. Neither corresponds
to functional importance for BLV navigation or interaction.

\textbf{VQA-based assistance:} Answering specific questions (``Is the
crosswalk light green?'') requires the user to know what to ask. Open-ended
scene description is essential for unknown environments.

\textbf{Prior assistive captioning:} Work on VizWiz-focused models
improves answer accuracy for BLV-generated images but does not address
the ordering of information in free-form descriptions.

The key missing component is a notion of \emph{salience} grounded in
human-centered priorities rather than photographic composition.

\section{Core Mechanism}

The Salience-LLaVA framework has two components:

\textbf{Salience scoring module:} A lightweight MLP head added to the
VLM's visual encoder that predicts a scalar salience score $s_i \in [0,1]$
for each detected bounding box $o_i$. The module is trained with a
point-wise ranking loss against human-annotated salience ranks.

\textbf{Salience-conditioned caption generation:} The caption generator
receives the ranked object list $[o_{\pi(1)}, o_{\pi(2)}, \ldots]$
(sorted by predicted salience) as soft prompt tokens prepended to the
image features. The LM is then trained to generate captions that mention
objects in the order they appear in the ranked list.

At inference, salience scoring and caption generation run in sequence:
score all detected objects, sort by salience, generate a caption
conditioned on the sorted list.

\section{Technical Formulation}

Let $O = \{o_1, \ldots, o_n\}$ be the detected objects with ground-truth
salience scores $s^*_i$. The salience head predicts $\hat{s}_i$ and
the ranking loss penalizes misordered pairs:
\[
  \mathcal{L}_\text{rank} = \sum_{i,j:\, s^*_i > s^*_j}
    \max\!\left(0,\, \hat{s}_j - \hat{s}_i + \delta\right)
\]
with margin $\delta = 0.1$.

The total training loss is:
\[
  \mathcal{L} = \mathcal{L}_\text{CE}(y, \hat{y}) + \lambda \mathcal{L}_\text{rank}
\]
where $\mathcal{L}_\text{CE}$ is cross-entropy on the caption sequence
$y$ and $\lambda = 0.3$.

The implicit ordering $\hat{\pi}$ extracted from the generated caption is
defined by the first mention position of each object:
\[
  \hat{\pi}(i) = \min\!\left\{t : o_i \text{ is mentioned at token position } t\right\}
\]
The ranking loss targets the ordering induced by $\hat{s}$ matching $s^*$.

\section{Learning or Inference Procedure}

\textbf{Dataset construction:}
\begin{enumerate}[itemsep=2pt,parsep=0pt]
  \item \textbf{Salience COCO:} 30k MSCOCO images; annotators rank all
    detected objects by importance for a BLV person navigating the scene.
  \item \textbf{Salience Flickr:} 10k Flickr images with the same protocol.
  \item \textbf{Salience VizWiz:} 15k images from BLV users' own
    photographs; salience ranks derived from crowdsourced question answers.
\end{enumerate}

\textbf{Training:} Fine-tune LLaVA-1.6 on Salience COCO and Salience
Flickr for 3 epochs with $\lambda = 0.3$; evaluate zero-shot transfer
to Salience VizWiz.

\textbf{Inference:} Run object detection (DINO-v2), score salience,
sort, generate caption with sorted soft prompt. Total overhead: $+38$ ms
per image on an A100.

\section{What the Guarantee Says}

The framework guarantees salience ordering in expectation over the
training distribution: at convergence, the expected rank of the
highest-salience object's first mention position approaches 1.
The guarantee is probabilistic—individual captions may reorder objects
due to LM fluency constraints—but the aggregate improvement is
statistically significant ($p < 0.001$).

\section{Experimental Findings}

\begin{itemize}[itemsep=2pt,parsep=0pt]
  \item \textbf{Top-1 salience ordering:} Salience-LLaVA 84.3\% vs.\
    LLaVA-1.6 51.2\%, GPT-4V-captioning 58.7\%.
  \item \textbf{VizWiz transfer:} Salience ordering F1 $+12.3$ pp over
    LLaVA-1.6 without VizWiz fine-tuning.
  \item \textbf{Caption quality:} SPICE +2.1, CIDEr +5.3 over LLaVA-1.6,
    confirming salience ordering does not hurt coverage.
  \item \textbf{User study (n=24 BLV participants):} Helpfulness
    4.2/5 vs.\ 3.1/5 for LLaVA; 20/24 preferred Salience-LLaVA for
    navigation tasks; most cited ``important info came first.''
\end{itemize}

\section{Ablations and Interpretation}

\textbf{Soft prompt vs.\ hard ordering:} Providing the ranked object list
as a hard prefix (forcing the caption to start with each object in order)
achieves 100\% top-1 ordering but reduces fluency (SPICE $-8.4$). Soft
prompting achieves 84.3\% ordering with only $+2.1$ SPICE gain.

\textbf{Salience source:} Using ground-truth salience ranks at inference
yields 92.1\% top-1 ordering, establishing an upper bound; the predicted
salience scores are 93\% rank-correlated with human rankings.

\textbf{$\lambda$ ablation:} $\lambda = 0$ (no ranking loss) reduces
top-1 ordering to 54.1\%; $\lambda = 1.0$ improves ordering to 87.2\%
but reduces CIDEr by 11 points; $\lambda = 0.3$ balances both.

\textbf{Domain specificity:} Salience annotations vary across domains:
in traffic scenes, moving objects (cars, cyclists) are highest salience;
in indoor scenes, social agents (people) dominate. The model generalizes
well within COCO categories but may require domain-specific annotation
for specialized environments.

\section*{Reference}

Jiazhao Liang, Hao Huang, Shuaihang Yuan, et al.
\textbf{Focus Where It Counts: A Salience-Driven Vision-Language
Model for Low Vision Assistance.}
arXiv:2608.28218, ECCV 2026.
\url{https://arxiv.org/abs/2608.28218}

\end{document}
